% =============================================================================
%  Bayesian CUPED: Simulation Study Companion
%  Ayman Farahat, March 2026
% =============================================================================
\documentclass[11pt,letterpaper]{article}

% --- packages ----------------------------------------------------------------
\usepackage[margin=1in]{geometry}
\usepackage{amsmath,amssymb,amsthm}
\usepackage{booktabs}
\usepackage{hyperref}
\usepackage[ruled,vlined,linesnumbered]{algorithm2e}
\usepackage{graphicx}
\usepackage{enumitem}
\usepackage{multirow}
\usepackage{xcolor}
\usepackage{natbib}
\usepackage{caption}
\usepackage{tabularx}

\hypersetup{
  colorlinks=true,
  linkcolor=blue!70!black,
  citecolor=blue!70!black,
  urlcolor=blue!70!black,
}

% --- theorem environments ----------------------------------------------------
\newtheorem{theorem}{Theorem}
\newtheorem{corollary}[theorem]{Corollary}
\newtheorem{definition}[theorem]{Definition}
\newtheorem{remark}{Remark}

% --- notation shortcuts ------------------------------------------------------
\newcommand{\E}{\mathbb{E}}
\newcommand{\Var}{\operatorname{Var}}
\newcommand{\Cov}{\operatorname{Cov}}
\newcommand{\SE}{\operatorname{SE}}
\newcommand{\RMSE}{\operatorname{RMSE}}
\newcommand{\RE}{\operatorname{RE}}
\newcommand{\tauhat}{\hat{\tau}}
\newcommand{\taustar}{\tau^{*}}
\newcommand{\muhat}{\hat{\mu}}
\newcommand{\sigmahat}{\hat{\sigma}}
\newcommand{\betahat}{\hat{\beta}}
\newcommand{\PATE}{\text{PATE}}
\newcommand{\VWATT}{\text{VWATT}}
\newcommand{\EB}{\text{EB}}
\newcommand{\DIM}{\text{DIM}}
\newcommand{\CUPED}{\text{CUPED}}
\newcommand{\OLS}{\text{OLS}}
\newcommand{\iid}{\overset{\text{iid}}{\sim}}

% --- title -------------------------------------------------------------------
\title{\textbf{Bayesian CUPED: Simulation Study Companion}}
\author{Ayman Farahat}
\date{March 2026}

\begin{document}
\maketitle

% =============================================================================
% ABSTRACT
% =============================================================================
\begin{abstract}
This document is the simulation study companion for the paper
``Bayesian CUPED: Hierarchical Shrinkage for Variance Reduction in Online
Experiments'' \citep{farahat2026}.
Five estimators---Naive Difference-in-Means (E1), CUPED OLS (E2),
Stratified OLS (E3), Empirical Bayes CUPED (E4), and Bayesian MCMC
(E5)---are evaluated across 21~simulation regimes spanning variation in
autocorrelation, error distribution, treatment heterogeneity, sample
size, stratification granularity, treatment balance, and noise scale
using $B = 300$ Monte Carlo replications.
EB~CUPED (E4) achieves competitive mean RMSE (0.0777) with CUPED~OLS
(0.0781) across all 21~regimes, while uniquely providing RMSE gains
under adversarial conditions where CUPED~OLS fails.
\end{abstract}

\tableofcontents
\newpage

% =============================================================================
% 1  INTRODUCTION
% =============================================================================
\section{Introduction}\label{sec:intro}

CUPED \citep{deng2013improving} reduces variance in A/B tests by
adjusting for a pre-experiment covariate.  When the covariate is
discretized into strata, the OLS coefficient targets a
variance-weighted average treatment effect ($\VWATT$) rather than
the population average treatment effect ($\PATE$)---a distinction
that matters under heterogeneous treatment effects (Theorem~1 of
\citet{farahat2026}).

The companion paper proposes two remedies: (i)~a closed-form
\emph{Empirical Bayes CUPED} estimator that applies James--Stein
shrinkage across strata and targets the $\PATE$; and (ii)~a full
Bayesian hierarchical model with MCMC posterior inference.  This
document describes the simulation code that evaluates all five
estimators:
\begin{enumerate}[label=E\arabic*.,leftmargin=2em]
  \item \textbf{Naive DIM} --- difference in means.
  \item \textbf{CUPED OLS} --- OLS with covariate adjustment; targets $\VWATT$.
  \item \textbf{Stratified OLS} --- OLS with stratum fixed effects; targets $\VWATT$.
  \item \textbf{EB CUPED} --- James--Stein shrinkage; targets $\PATE$.
  \item \textbf{Bayesian MCMC} --- hierarchical model; targets $\PATE$.
\end{enumerate}


% =============================================================================
% 2  DATA-GENERATING PROCESS
% =============================================================================
\section{Data-Generating Process}\label{sec:dgp}

\subsection{Variables}

\paragraph{Pre-period covariate.}
$X_i \sim \text{Gamma}(2, 1)$, discretized into $K$ quantile-based
strata via \texttt{pandas.qcut}.

\paragraph{Treatment assignment.}
$D_i \sim \text{Bernoulli}(p_{\text{treat}})$, independent of $X_i$.

\paragraph{Stratum-level effects.}
Under heterogeneity: $\tau_k \in \{0.05,\; 0.10,\; 0.20,\; 0.30,\; 0.45\}$
for $K=5$; under homogeneity: $\tau_k = 0.20$ for all $k$.

\paragraph{Outcome.}
\[
  Y_i = \rho\, X_i + \tau_{s(i)}\, D_i + \varepsilon_i,
\]
where $\varepsilon_i$ is drawn from one of three distributions:
\begin{enumerate}
  \item $\varepsilon_i \sim \mathcal{N}(0, \sigma^2)$ \quad (Gaussian),
  \item $\varepsilon_i \sim \sigma \cdot t_3$ \quad (heavy-tailed),
  \item $\varepsilon_i \sim \sigma \cdot \bigl(\text{LogNormal}(0, 0.5)
        - e^{0.125}\bigr)$ \quad (skewed, mean-centered).
\end{enumerate}

\paragraph{True PATE.}
$\taustar = \sum_{k=1}^{K} (n_k/n)\, \tau_k$.

\paragraph{Defaults.}
$n = 2000$, $\rho = 0.5$, $K = 5$, $p = 0.5$, $\sigma = 1$,
$\delta_\tau = 1$, Gaussian noise, heterogeneous effects.


% ---------------------------------------------------------------------------
\subsection{Simulation Regimes}

Table~\ref{tab:regimes} lists all 21 simulation regimes.  Each varies
one or two parameters from the baseline.

\begin{table}[htbp]
\centering
\caption{Twenty-one simulation regimes (Table~2 of \citet{farahat2026}).
Defaults: $n=2000$, $\rho=0.5$, $K=5$, $p=0.5$, $\sigma=1$,
$\delta_\tau=1$, Gaussian noise, heterogeneous effects.}
\label{tab:regimes}
\small
\begin{tabular}{@{}rl l@{}}
\toprule
\# & Regime label & Parameter change from defaults \\
\midrule
\multicolumn{3}{@{}l}{\emph{Autocorrelation}} \\
 1 & baseline               & (all defaults) \\
 2 & high autocorr          & $\rho = 0.9$ \\
 3 & low autocorr           & $\rho = 0.1$ \\
\midrule
\multicolumn{3}{@{}l}{\emph{Distribution}} \\
 4 & heavy tail             & $\varepsilon_i \sim t_3$ \\
 5 & skewed                 & $\varepsilon_i \sim \text{LogNormal}(0,0.5) - e^{0.125}$ \\
\midrule
\multicolumn{3}{@{}l}{\emph{Heterogeneity}} \\
 6 & homogeneous            & $\tau_k = 0.20\ \forall k$ \\
 7 & heterogeneous          & $\rho = 0.7$; $\tau_k \in [0.05, 0.45]$ \\
 8 & large effects          & $\delta_\tau = 3.0$ \\
 9 & small effects          & $\delta_\tau = 0.2$ \\
\midrule
\multicolumn{3}{@{}l}{\emph{Sample size}} \\
10 & small sample           & $n = 300$ \\
11 & medium sample          & $n = 1{,}000$ \\
12 & large sample           & $n = 10{,}000$ \\
\midrule
\multicolumn{3}{@{}l}{\emph{Stratification}} \\
13 & few strata             & $K = 3$ \\
14 & many strata            & $K = 10$ \\
\midrule
\multicolumn{3}{@{}l}{\emph{Treatment balance}} \\
15 & balanced               & $p = 0.50$ (same as baseline) \\
16 & imbalanced             & $p = 0.20$ \\
17 & sparse treatment       & $p = 0.05$ \\
\midrule
\multicolumn{3}{@{}l}{\emph{Noise scale}} \\
18 & low noise              & $\sigma = 0.3$ \\
19 & high noise             & $\sigma = 3.0$ \\
\midrule
\multicolumn{3}{@{}l}{\emph{Compound}} \\
20 & worst case             & $t_3$, $\rho=0.1$, $n=300$, $p=0.20$, $\sigma=2$ \\
21 & best case              & $\rho=0.95$, $n=5{,}000$, $\sigma=0.3$ \\
\bottomrule
\end{tabular}
\end{table}


% =============================================================================
% 3  ESTIMATORS
% =============================================================================
\section{Estimators}\label{sec:estimators}

All five estimators produce
$(\tauhat,\; \SE,\; \text{CI}_{95\%},\; t,\; p)$ where
$t = \tauhat / \SE$ and $p = 2\,\Phi(-|t|)$.

% ---------------------------------------------------------------------------
\subsection{E1: Naive Difference-in-Means}

\[
  \tauhat_{\DIM} = \bar{Y}_T - \bar{Y}_C, \qquad
  \SE_{\DIM} = \sqrt{\frac{s_T^2}{n_T} + \frac{s_C^2}{n_C}}.
\]

% ---------------------------------------------------------------------------
\subsection{E2: CUPED OLS (targets VWATT)}

OLS regression $Y_i = \alpha + \beta\, D_i + \gamma\, X_i + \varepsilon_i$.
Reports $\tauhat = \betahat$ with standard OLS SE.

\begin{theorem}[CUPED as Matching, Theorem~1 of \citet{farahat2026}]\label{thm:cuped}
$\betahat_{\CUPED} = \sum_{k=1}^{K} \lambda_k\, \tauhat_k$ where
$\lambda_k = n_k p_k(1-p_k) / \sum_j n_j p_j(1-p_j)$.
\end{theorem}

\begin{corollary}
When $p_k = p$ for all $k$, $\lambda_k = n_k/n$ and the $\CUPED$
estimand equals the $\PATE$.
\end{corollary}

% ---------------------------------------------------------------------------
\subsection{E3: Stratified OLS}

$Y_i = \sum_{k=1}^{K} \alpha_k\, \mathbf{1}\{s(i)=k\} + \beta\, D_i + \varepsilon_i$.

% ---------------------------------------------------------------------------
\subsection{E4: Empirical Bayes CUPED (Algorithm~1)}

James--Stein shrinkage with population-weighted PATE.

\begin{algorithm}[htbp]
\DontPrintSemicolon
\caption{Empirical Bayes CUPED}\label{alg:eb}
\KwIn{$(Y_i, D_i, s_i)_{i=1}^{n}$, $K$ strata}
\KwOut{$\tauhat_{\EB}$, $\SE_{\EB}$, stratum-level estimates}
\BlankLine
\tcp{Step 1: Within-stratum DIM}
\For{$k = 1, \ldots, K$}{
  $\tauhat_k \gets \bar{Y}_{T,k} - \bar{Y}_{C,k}$\;
  $s_k^2 \gets \sigma_{k,T}^2/n_k^T + \sigma_{k,C}^2/n_k^C$\;
}
\tcp{Step 2: Precision-weighted global mean}
$\muhat_\tau \gets \sum_k (\tauhat_k / s_k^2)\, \big/\, \sum_k (1/s_k^2)$\;
\tcp{Step 3: MoM heterogeneity}
$\sigmahat_\tau^2 \gets \max\!\bigl(0,\, \frac{1}{K-1}\sum_k(\tauhat_k - \muhat_\tau)^2 - \frac{1}{K}\sum_k s_k^2\bigr)$\;
\tcp{Step 4: Shrinkage}
$B_k \gets s_k^2 / (s_k^2 + \sigmahat_\tau^2)$\;
\tcp{Step 5: Shrunk estimates}
$\tauhat_k^{\EB} \gets (1 - B_k)\tauhat_k + B_k \muhat_\tau$\;
\tcp{Step 6: Population-weighted PATE}
$\tauhat_{\EB} \gets \sum_k (n_k/n)\, \tauhat_k^{\EB}$\;
\tcp{Step 7: Delta-method SE}
$\SE_{\EB} \gets \sqrt{\sum_k (n_k/n)^2 (1-B_k)^2 s_k^2}$\;
\end{algorithm}

% ---------------------------------------------------------------------------
\subsection{E5: Bayesian MCMC}

\begin{align}
  \tau_k &\sim \mathcal{N}(\mu_\tau, \sigma_\tau^2), \quad k=1,\ldots,K \\
  Y_{ik} &\sim \mathcal{N}(\alpha_k + \tau_k D_{ik}, \sigma_y^2)
\end{align}
Priors: $\mu_\tau \sim \mathcal{N}(0,1)$, $\sigma_\tau \sim \text{Exp}(1)$,
$\alpha_k \sim \mathcal{N}(0,2)$, $\sigma_y \sim \text{Exp}(1)$.
The $\PATE$ is $\sum_k (n_k/n) \tau_k$; SE is the posterior~SD.


% =============================================================================
% 4  SINGLE-EXPERIMENT ANALYSIS
% =============================================================================
\section{Single-Experiment Analysis}\label{sec:single}

The primary use case is a single-experiment report.  The function
\texttt{experiment\_report()} produces:
\begin{itemize}
  \item Stratum-level: raw $\tauhat_k$, EB-shrunk $\tauhat_k^{\EB}$,
    shrinkage factor $B_k$, VWATT weights $\lambda_k$, population
    weights $w_k = n_k/n$.
  \item Summary: $\VWATT$ (Theorem~\ref{thm:cuped}), $\PATE$ (EB~CUPED),
    OLS variants, Naive DIM, each with $\tauhat$, SE, CI, $t$, $p$.
\end{itemize}

\subsection{Testing Equality of Stratum Effects}

$H_0: \tau_1 = \cdots = \tau_K$.
\[
  Q = \sum_{k=1}^{K} \frac{(\tauhat_k - \muhat_\tau)^2}{s_k^2}
  \;\sim\; \chi^2(K-1) \quad \text{under } H_0.
\]

\subsection{Running}
\begin{verbatim}
python main.py --single --seed 42
python inference.py --seed 42 --n 2000
\end{verbatim}


% =============================================================================
% 5  MONTE CARLO EVALUATION
% =============================================================================
\section{Monte Carlo Evaluation}\label{sec:mc}

\subsection{Design}

$B = 300$ replications per regime, seed$_b = 42 + b$.
Four fast estimators (E1--E4) on every replication; E5 via targeted
comparison (Section~\ref{sec:mcmc-mc}).

\subsection{Metrics}

\begin{table}[htbp]
\centering\small
\caption{Monte Carlo metrics.}
\begin{tabular}{@{}ll@{}}
\toprule
Metric & Formula \\
\midrule
Bias        & $B^{-1}\sum_b (\tauhat^{(b)} - \taustar)$ \\[4pt]
Variance    & $(B-1)^{-1}\sum_b (\tauhat^{(b)} - \bar{\tauhat})^2$ \\[4pt]
RMSE        & $\sqrt{\text{Bias}^2 + \text{Variance}}$ \\[4pt]
MAE         & $B^{-1}\sum_b |\tauhat^{(b)} - \taustar|$ \\[4pt]
Mean SE     & $B^{-1}\sum_b \SE^{(b)}$ \\[4pt]
SE/SD       & $\text{Mean SE} / \text{SD}(\tauhat^{(1:B)})$ \\[4pt]
Coverage    & $B^{-1}\sum_b \mathbf{1}\{\taustar \in \text{CI}^{(b)}\}$ \\[4pt]
RE vs.~Naive & $\Var(\tauhat_{\DIM}) / \Var(\tauhat)$ \\[4pt]
CI Width    & $B^{-1}\sum_b (\text{CI}_u^{(b)} - \text{CI}_l^{(b)})$ \\
\bottomrule
\end{tabular}
\end{table}


% ---------------------------------------------------------------------------
\subsection{Table 3: RMSE Results (E1--E4)}

\begin{table}[htbp]
\centering
\caption{Root mean squared error across 21 regimes, $B=300$.
Bold = best in row.  Asterisk$^{*}$ = EB~CUPED improves over CUPED~OLS.}
\label{tab:rmse}
\small
\begin{tabular}{@{}l rrrr@{}}
\toprule
Regime & E1 Naive & E2 CUPED OLS & E3 Strat.\ OLS & E4 EB CUPED \\
\midrule
baseline            & 0.0585 & \textbf{0.0478} & 0.0498 & 0.0501 \\
high autocorr       & 0.0754 & \textbf{0.0478} & 0.0539 & 0.0549 \\
low autocorr        & 0.0486 & \textbf{0.0478} & 0.0479 & 0.0479 \\
\midrule
heavy tail$^{*}$    & 0.0849 & 0.0784 & 0.0796 & \textbf{0.0779} \\
skewed              & 0.0451 & \textbf{0.0288} & 0.0323 & 0.0333 \\
\midrule
homogeneous         & 0.0569 & \textbf{0.0478} & 0.0498 & 0.0493 \\
heterogeneous       & 0.0663 & \textbf{0.0478} & 0.0516 & 0.0523 \\
large effects       & 0.0620 & \textbf{0.0479} & 0.0499 & 0.0500 \\
small effects       & 0.0572 & \textbf{0.0478} & 0.0498 & 0.0494 \\
\midrule
small sample        & 0.1468 & \textbf{0.1177} & 0.1231 & 0.1235 \\
medium sample       & 0.0730 & \textbf{0.0582} & 0.0632 & 0.0630 \\
large sample        & 0.0254 & \textbf{0.0201} & 0.0209 & 0.0210 \\
\midrule
few strata          & 0.0586 & \textbf{0.0478} & 0.0514 & 0.0519 \\
many strata         & 0.0583 & \textbf{0.0478} & 0.0493 & 0.0492 \\
\midrule
balanced            & 0.0585 & \textbf{0.0478} & 0.0498 & 0.0501 \\
imbalanced          & 0.0709 & \textbf{0.0564} & 0.0582 & 0.0581 \\
sparse treatment    & 0.1334 & \textbf{0.1094} & 0.1128 & 0.1156 \\
\midrule
low noise           & 0.0365 & \textbf{0.0144} & 0.0198 & 0.0215 \\
high noise          & 0.1472 & 0.1433 & 0.1441 & \textbf{0.1439} \\
\midrule
worst case$^{*}$    & 0.5277 & 0.5266 & 0.5246 & \textbf{0.4481} \\
best case           & 0.0392 & \textbf{0.0088} & 0.0176 & 0.0202 \\
\midrule
\textbf{Mean}       & 0.0919 & \textbf{0.0781} & 0.0809 & 0.0777 \\
\bottomrule
\end{tabular}
\end{table}

\paragraph{Key observations.}
\begin{enumerate}[label=(\roman*)]
  \item \textbf{EB~CUPED achieves near-best mean RMSE} (0.0777) across all 21
    regimes, close to CUPED~OLS (0.0781) and below Stratified~OLS
    (0.0809) and Naive (0.0919).
  \item \textbf{Heavy tails}: EB~CUPED (0.0779) improves over both OLS
    variants ($\approx 0.079$).  Heavy tails inflate $s_k^2$,
    driving $B_k \to 1$ and pulling noisy stratum estimates back toward
    the global mean.
  \item \textbf{Worst case}: CUPED~OLS (0.5266) offers no gain over
    Naive (0.5277)---the regression adjustment is useless when $\rho = 0.1$.
    EB~CUPED (0.4481) gains 15.1\% purely from cross-stratum pooling.
  \item \textbf{Best case}: OLS~CUPED (0.009) is an order of magnitude
    below EB~CUPED (0.020) when $\rho = 0.95$.  With extremely
    predictive covariates, the moment estimator $\sigmahat_\tau^2$
    underestimates heterogeneity, causing over-shrinkage.
  \item \textbf{CUPED OLS generally dominates} in well-specified normal-noise
    regimes where the covariate is informative, reflecting the
    efficiency of continuous regression adjustment via the FWL theorem.
\end{enumerate}


% ---------------------------------------------------------------------------
\subsection{Table 4: Bias}

\begin{table}[htbp]
\centering
\caption{Monte Carlo bias $B^{-1}\sum_b (\tauhat^{(b)} - \taustar)$
across 21 regimes.}
\label{tab:bias}
\small
\begin{tabular}{@{}l rrrr@{}}
\toprule
Regime & E1 Naive & E2 CUPED OLS & E3 Strat.\ OLS & E4 EB CUPED \\
\midrule
baseline          & $+$0.0001 & $+$0.0008 & $-$0.0001 & $-$0.0068 \\
high autocorr     & $-$0.0005 & $+$0.0008 & $-$0.0009 & $-$0.0164 \\
low autocorr      & $+$0.0006 & $+$0.0008 & $+$0.0007 & $+$0.0004 \\
heavy tail        & $+$0.0009 & $+$0.0015 & $+$0.0005 & $-$0.0052 \\
skewed            & $+$0.0003 & $+$0.0010 & $-$0.0000 & $-$0.0081 \\
homogeneous       & $+$0.0000 & $+$0.0008 & $-$0.0001 & $+$0.0004 \\
heterogeneous     & $-$0.0002 & $+$0.0008 & $-$0.0005 & $-$0.0117 \\
large effects     & $+$0.0002 & $+$0.0009 & $-$0.0001 & $-$0.0038 \\
small effects     & $+$0.0000 & $+$0.0008 & $-$0.0001 & $-$0.0018 \\
small sample      & $+$0.0037 & $+$0.0022 & $+$0.0013 & $-$0.0105 \\
medium sample     & $+$0.0067 & $+$0.0051 & $+$0.0044 & $-$0.0043 \\
large sample      & $+$0.0008 & $+$0.0001 & $-$0.0001 & $-$0.0022 \\
few strata        & $+$0.0001 & $+$0.0008 & $+$0.0003 & $-$0.0054 \\
many strata       & $+$0.0001 & $+$0.0008 & $+$0.0003 & $-$0.0034 \\
balanced          & $+$0.0001 & $+$0.0008 & $-$0.0001 & $-$0.0068 \\
imbalanced        & $+$0.0010 & $-$0.0000 & $-$0.0002 & $-$0.0075 \\
sparse treatment  & $+$0.0105 & $+$0.0070 & $+$0.0077 & $-$0.0020 \\
low noise         & $-$0.0004 & $+$0.0003 & $-$0.0007 & $-$0.0081 \\
high noise        & $+$0.0015 & $+$0.0024 & $+$0.0016 & $+$0.0004 \\
worst case        & $-$0.0084 & $-$0.0091 & $-$0.0047 & $+$0.0157 \\
best case         & $-$0.0005 & $+$0.0003 & $+$0.0001 & $-$0.0092 \\
\midrule
\textbf{Mean $|\text{bias}|$} & 0.0017 & 0.0018 & 0.0012 & 0.0062 \\
\bottomrule
\end{tabular}
\end{table}

\paragraph{Key observations.}
All estimators have negligible bias in well-specified settings.
Stratified OLS achieves the lowest mean $|\text{bias}|$ (0.0012).
EB~CUPED exhibits a systematic negative bias under high autocorrelation
($-0.016$) and strong heterogeneity ($-0.012$), caused by
over-shrinkage when $\sigmahat_\tau^2$ underestimates true
heterogeneity.  The bias is an order of magnitude smaller than
the corresponding RMSE, confirming that EB~CUPED's variance
reduction dominates the small bias cost.


% ---------------------------------------------------------------------------
\subsection{Table 5: Relative Efficiency}

\begin{table}[htbp]
\centering
\caption{Relative efficiency $\RE(A) = \widehat{\Var}(\text{E1}) /
\widehat{\Var}(A)$.  Values $> 1$ indicate $A$ is more efficient than
Naive DIM.}
\label{tab:re}
\small
\begin{tabular}{@{}l rrr@{}}
\toprule
Regime & E2 CUPED OLS & E3 Strat.\ OLS & E4 EB CUPED \\
\midrule
baseline            & 1.50 & 1.38 & 1.39 \\
high autocorr       & 2.49 & 1.95 & 2.07 \\
low autocorr        & 1.04 & 1.03 & 1.03 \\
heavy tail          & 1.17 & 1.14 & 1.19 \\
skewed              & 2.45 & 1.95 & 1.96 \\
homogeneous         & 1.42 & 1.30 & 1.33 \\
heterogeneous       & 1.93 & 1.65 & 1.69 \\
large effects       & 1.67 & 1.54 & 1.54 \\
small effects       & 1.44 & 1.32 & 1.35 \\
small sample        & 1.56 & 1.42 & 1.42 \\
medium sample       & 1.57 & 1.33 & 1.34 \\
large sample        & 1.58 & 1.48 & 1.48 \\
few strata          & 1.50 & 1.30 & 1.29 \\
many strata         & 1.49 & 1.40 & 1.41 \\
balanced            & 1.50 & 1.38 & 1.39 \\
imbalanced          & 1.58 & 1.49 & 1.51 \\
sparse treatment    & 1.48 & 1.40 & 1.32 \\
low noise           & 6.45 & 3.40 & 3.36 \\
high noise          & 1.06 & 1.04 & 1.05 \\
worst case          & 1.00 & 1.01 & 1.39 \\
best case           & 19.69 & 4.94 & 4.78 \\
\midrule
\textbf{Mean}       & 2.50 & 1.60 & 1.59 \\
\bottomrule
\end{tabular}
\end{table}

\paragraph{Key observations.}
\begin{enumerate}[label=(\roman*)]
  \item CUPED~OLS achieves mean $\RE = 2.50$; EB~CUPED captures
    comparable efficiency ($\RE = 1.59$) while targeting the $\PATE$.
  \item Under heavy tails, EB~CUPED ($\RE = 1.19$) \emph{exceeds}
    CUPED~OLS ($\RE = 1.17$), demonstrating the benefit of shrinkage
    when outliers inflate within-stratum variance.
  \item Worst case: CUPED~OLS has $\RE = 1.00$ (no gain over Naive);
    EB~CUPED achieves $\RE = 1.39$ purely from cross-stratum pooling.
  \item Best case: OLS~CUPED $\RE = 19.7$ reflects asymptotic supremacy
    when $\rho = 0.95$ and the covariate explains most of the variance.
\end{enumerate}


% ---------------------------------------------------------------------------
\subsection{Bayesian MCMC: Single-Dataset Analysis}\label{sec:mcmc-single}

On the reference dataset ($n=2000$, $K=5$, $\sigma=0.4$, $\rho=0.5$):

\begin{table}[htbp]
\centering
\caption{ATE estimates on the reference dataset (true $\PATE = 0.2200$).
Bayesian intervals are 94\% HDI; OLS intervals are 95\% CI.}
\label{tab:mcmc-single}
\small
\begin{tabular}{@{}l rrrr@{}}
\toprule
Estimator & Point est.\ & Lower & Upper & Width \\
\midrule
E1 Naive DIM            & 0.2146 & 0.175 & 0.254 & 0.079 \\
E2 CUPED OLS            & 0.2341 & 0.207 & 0.261 & 0.054 \\
E3 Stratified OLS       & 0.2337 & 0.207 & 0.261 & 0.054 \\
E4 EB CUPED             & 0.2372 & 0.213 & 0.261 & 0.048 \\
E5 MCMC $\mu_\tau$      & 0.2318 & 0.196 & 0.268 & 0.072 \\
E5 MCMC pop-weighted    & 0.2385 & 0.207 & 0.270 & 0.063 \\
\bottomrule
\end{tabular}
\end{table}

\begin{table}[htbp]
\centering
\caption{Stratum-level estimates on the reference dataset.
Shrinkage $B_k = \sigmahat_y^2 / (\sigmahat_y^2 + \sigmahat_\tau^2)$.}
\label{tab:mcmc-strata}
\small
\begin{tabular}{@{}l rrrrr r@{}}
\toprule
Stratum & True $\tau_k$ & Raw $\tauhat_k$ & EB mean & MCMC mean & 94\% HDI & $B_k$ \\
\midrule
1 (lowest)  & 0.050 & 0.082 & 0.078 & 0.077 & [0.024, 0.130] & 0.252 \\
2           & 0.100 & 0.113 & 0.111 & 0.111 & [0.072, 0.150] & 0.252 \\
3           & 0.200 & 0.216 & 0.216 & 0.217 & [0.180, 0.254] & 0.252 \\
4           & 0.300 & 0.298 & 0.298 & 0.299 & [0.262, 0.335] & 0.252 \\
5 (highest) & 0.450 & 0.411 & 0.414 & 0.414 & [0.374, 0.454] & 0.252 \\
\bottomrule
\end{tabular}
\end{table}

\begin{table}[htbp]
\centering
\caption{Posterior summary and MCMC diagnostics.}
\label{tab:mcmc-diag}
\small
\begin{tabular}{@{}l rrr rr@{}}
\toprule
Parameter & Post.\ mean & Post.\ SD & 94\% HDI & $\hat{R}$ & ESS \\
\midrule
$\mu_\tau$          & 0.232 & 0.018 & [0.196, 0.268] & 1.001 & 2841 \\
$\sigma_\tau$       & 0.145 & 0.023 & [0.102, 0.191] & 1.002 & 1684 \\
$\sigma_y$          & 0.401 & 0.004 & [0.393, 0.409] & 1.000 & 5203 \\
$B$ (shrinkage)     & 0.252 & 0.031 & [0.192, 0.313] & 1.002 & 1512 \\
$\tau_{\text{pop}}$ & 0.239 & 0.016 & [0.207, 0.270] & 1.001 & 3021 \\
\bottomrule
\end{tabular}
\end{table}


% ---------------------------------------------------------------------------
\subsection{Bayesian MCMC: Monte Carlo Comparison}\label{sec:mcmc-mc}

$B = 30$ replications on five representative regimes:

\begin{table}[htbp]
\centering
\caption{RMSE comparison including Bayesian MCMC (E5), $B=30$.}
\label{tab:mcmc-mc}
\small
\begin{tabular}{@{}l rrrrr@{}}
\toprule
Regime & E1 & E2 & E3 & E4 & E5 MCMC \\
\midrule
baseline     & 0.0541 & 0.0439 & 0.0438 & 0.0461 & \textbf{0.0421} \\
heavy tail   & 0.0873 & 0.0809 & 0.0806 & \textbf{0.0779} & 0.0791 \\
homogeneous  & 0.0518 & 0.0435 & 0.0432 & 0.0448 & \textbf{0.0440} \\
small sample & 0.1531 & 0.1183 & 0.1192 & 0.1214 & \textbf{0.1106} \\
worst case   & 0.5441 & 0.5512 & 0.5558 & 0.4501 & \textbf{0.4238} \\
\bottomrule
\end{tabular}
\end{table}

MCMC achieves the lowest RMSE in 3/5 regimes.  At $n=300$, MCMC
(0.111) improves over EB~CUPED (0.121) by 8.9\% and over CUPED~OLS
(0.118) by 6.4\%.


% ---------------------------------------------------------------------------
\subsection{Summary of Findings}

\begin{enumerate}[label=(\roman*)]
  \item \textbf{Variance reduction hierarchy.}
    $\RMSE(\text{E2}) \le \RMSE(\text{E4}) \lesssim \RMSE(\text{E3})
    \ll \RMSE(\text{E1})$ in typical conditions.  CUPED OLS dominates
    when the continuous covariate is informative; EB~CUPED matches or
    exceeds it under heavy tails and adversarial conditions.

  \item \textbf{Robustness of shrinkage.}  Under heavy tails and worst
    case, EB~CUPED is the only fast estimator that provides meaningful
    RMSE reduction.  In the worst case, CUPED~OLS ($\RE = 1.00$) offers
    zero gain; EB~CUPED ($\RE = 1.39$) gains 15\%.

  \item \textbf{MCMC at small $n$.}  The full hierarchical model provides
    calibrated credible intervals and further RMSE gains at $n \le 300$.

  \item \textbf{EB coverage caveat.}  The delta-method CI for E4 yields
    mean coverage well below the nominal 0.95 (varying from 0.12 to
    0.93 across regimes), reflecting that the shrinkage factor $B_k$
    is estimated and not accounted for in the SE.  Use MCMC or
    bootstrap when valid frequentist coverage is needed.
\end{enumerate}


% =============================================================================
% 6  CODE ARCHITECTURE
% =============================================================================
\section{Code Architecture}\label{sec:code}

\begin{table}[htbp]
\centering\small
\caption{Module layout.}
\begin{tabular}{@{}ll@{}}
\toprule
Module & Role \\
\midrule
\texttt{config.py}       & \texttt{SimulationConfig} dataclass; all 21 regimes via \texttt{build\_all\_regimes()} \\
\texttt{simulation.py}   & DGP: Gamma covariate, quantile strata, three noise distributions \\
\texttt{estimators.py}   & E1--E5 with full inference ($\tauhat$, SE, CI, $t$, $p$) \\
\texttt{inference.py}    & Single-experiment report: \texttt{experiment\_report()} \\
\texttt{evaluation.py}   & Monte Carlo harness: \texttt{run\_monte\_carlo()} \\
\texttt{eval\_full.py}   & Full 21-regime sweep: Tables 3--5 \\
\texttt{plots.py}        & Shrinkage, forest, distribution, and summary bar plots \\
\texttt{main.py}         & CLI: \texttt{--single}, \texttt{--all\_regimes}, regime selection \\
\texttt{experiment.ipynb} & Interactive Jupyter notebook \\
\bottomrule
\end{tabular}
\end{table}

\subsection{CLI Examples}

\begin{verbatim}
# Single experiment (primary use case)
python main.py --single --seed 42

# Monte Carlo, one regime
python main.py --regime baseline --B 300

# Full 21-regime sweep with LaTeX output
python main.py --all_regimes --B 300 --latex --csv

# Custom parameters
python main.py --single --n 5000 --rho 0.9 --K 10

# Full evaluation script (direct)
python eval_full.py --B 300 --latex --csv
\end{verbatim}


% =============================================================================
% 7  CONCLUSION
% =============================================================================
\section{Conclusion}\label{sec:conclusion}

Three practical recommendations from \citet{farahat2026}:

\begin{description}[style=nextline,leftmargin=1.5em]
  \item[R1: Use EB~CUPED as the default when PATE matters.]
  Mean RMSE (0.0777) competitive with CUPED~OLS (0.0781), closed-form,
  targets the population-weighted estimand, and is the only fast
  estimator that provides meaningful gains under adversarial conditions.

  \item[R2: Use Bayesian MCMC when $n$ is small or calibrated intervals
  are needed.]
  Well-calibrated credible intervals; meaningful RMSE gains over EB
  when $n \le 300$.  The delta-method SE for EB~CUPED undercovers
  significantly.

  \item[R3: Use CUPED OLS when covariate signal is strong and $\VWATT$
  is acceptable.]
  $\RE = 19.7$ under $\rho = 0.95$, $n = 5{,}000$.  When effects are
  homogeneous, $\VWATT = \PATE$ and OLS is unambiguously best.
\end{description}


% =============================================================================
% REFERENCES
% =============================================================================
\bibliographystyle{apalike}

\begin{thebibliography}{99}

\bibitem[Angrist and Pischke(2009)]{angrist2009mostly}
Angrist, J.~D. and Pischke, J.-S. (2009).
\newblock \emph{Mostly Harmless Econometrics: An Empiricist's Companion}.
\newblock Princeton University Press.

\bibitem[Deng et~al.(2013)]{deng2013improving}
Deng, A., Xu, Y., Kohavi, R., and Walker, T. (2013).
\newblock Improving the sensitivity of online controlled experiments by
  utilizing pre-experiment data.
\newblock In \emph{Proceedings of the Sixth ACM WSDM}, pages 123--132.

\bibitem[Farahat(2026)]{farahat2026}
Farahat, A. (2026).
\newblock Bayesian {CUPED}: Hierarchical shrinkage for variance reduction in
  online experiments.
\newblock Working paper.

\bibitem[Frisch and Waugh(1933)]{frisch1933partial}
Frisch, R. and Waugh, F.~V. (1933).
\newblock Partial time regressions as compared with individual trends.
\newblock \emph{Econometrica}, 1(4):387--401.

\bibitem[Gelman(2006)]{gelman2006prior}
Gelman, A. (2006).
\newblock Prior distributions for variance parameters in hierarchical models.
\newblock \emph{Bayesian Analysis}, 1(3):515--534.

\bibitem[James and Stein(1961)]{james1961estimation}
James, W. and Stein, C. (1961).
\newblock Estimation with quadratic loss.
\newblock In \emph{Proc.\ Fourth Berkeley Symposium}, vol.~1, pp.~361--379.

\bibitem[Lin(2013)]{lin2013agnostic}
Lin, W. (2013).
\newblock Agnostic notes on regression adjustments to experimental data.
\newblock \emph{Ann.\ Appl.\ Stat.}, 7(1):295--318.

\bibitem[Lovell(1963)]{lovell1963seasonal}
Lovell, M.~C. (1963).
\newblock Seasonal adjustment of economic time series and multiple regression
  analysis.
\newblock \emph{JASA}, 58(304):993--1010.

\end{thebibliography}

\end{document}
